\section{Truth and Grammar in TFL}

So far, despite some basic symbolization as stand-ins for sentences, our
discussion has been entirely informal: our arguments are built out of
ordinary English sentences; our symbolizations are likewise parphrases of
English.

This week we will take our first steps to \emph{formal logic} by describing
the \emph{syntax} of Truth-Functional Logic (TFL). We will learn what we mean
by truth-functional when we discuss the \emph{semantics} of TFL. For the time-
being we will concentrate on the \emph{form} in formal.

First and foremost, we want to know what ``counts'' as a sentence in TFL. We can
start with a few definitions about the things that might be sentences:
\emph{expressions}.

\begin{definition}[Basic TFL Expression]
  \label{def:tfl-basic-expression}
  A \emph{basic expression} in TFL is any of the following symbols:
  \begin{enumerate}
  \item atomic sentence letters: $A, B, C$
  \item connectives: \enot, \eand, \eor, \eif, \eiff
  \item brackets: (, )
  \end{enumerate}
  We will refer to the collection of basic expressions in TFL as its
  \emph{alphabet}.
\end{definition}

Note that the atomic sentences can be made both \emph{enumerable} and
\emph{potentially infintite} with the addition of subscripts and/or
superscripts: $A_4, B^2, C_1^2$, etc.

\begin{definition}[TFL Expression]
  \label{def:tfl-expression}
  An \emph{expression} in TFL is any sequence of basic TFL expressions.
\end{definition}

Many of these expressions will be nonsense. So we are really only interested in
some of these expressions: the \emph{sentences} of TFL. There are potentially
infinitely many of these, so we will have to define the sentences
\emph{recursively} using rules. That is, we will need to explain how to
construct sentences from expressions.

\begin{definition}[TFL Sentence]
  \label{def:tfl-sentence}
  A \emph{sentence} in TFL is defined by the following rules:
  \begin{enumerate}
  \item Every atomic sentence letter is a sentence.
  \item If \sphi is a sentence, then \enot\sphi is a sentence.
  \item If \sphi and \spsi are sentences, then (\ephi\eand\epsi) is a sentence.
  \item If \sphi and \spsi are sentences, then (\ephi\eor\epsi) is a sentence.
  \item If \sphi and \spsi are sentences, then (\ephi\eif\epsi) is a sentence.
  \item If \sphi and \spsi are sentences, then (\ephi\eiff\epsi) is a sentence.
  \item Nothing else is a sentence.
  \end{enumerate}
  A sentence is also called a \emph{well-formed formula} (\emph{wff}).
\end{definition}

This is called an \emph{inductive} defininition: it uses rules to build
infinitely many complex expressions from a small number of basic ones. Inductive
definitions form the backbone of most systems of deductive logic, both in syntax
and semantics. Notice that the last clause ensures that there are many TFL
expressions that are not TFL sentences. But of course every TFL sentence is a
TFL expression.

Before diving into more detailed issues of TFL syntax, we should pause to
comment on some of the differences between the connectives of TFL and their
english analogues.

First, TFL \snot is \emph{logical}. Specifically, it obeys three important
inter-related laws:

%\begin{argument}
%\item[\textbf{LEM:}] \ephi\eor\enot\ephi
%\item[\textbf{LNC:}] \enot(\ephi\eand\enot\ephi)
%\item[\textbf{DN:}]  \enot\enot\ephi\eiff\ephi  
%\end{argument}

yWe will explore these laws in more detail later. For now it suffices to say that
they mean that every sentence is either true or false, never anything in
between. And if A is true, then not-A is false. English ``not'' does not work
like this and can be applied to all sorts of English expressions, not just to
sentences.

Second, in English the order of two sentences connected by ``and'' is important.
We illustrated this last week with the following example:

\begin{argument}
\item[A:] John left.
\item[B:] Mary got upset.
\end{argument}

Suffice it to say, \emph{A and B} can express something different than
\emph{B and A} ijn English. Typically this will be something about the order
of the events described in $A$ and $B$, but also sometimes has to do with
causal relations between them. The TFL connective \eand, on the other hand,
is \emph{commutative}. $A\eand B$ always means the same thing as $B\eand A$.

Third, English ``or'' is often understood with an \emph{exclusive} reading:
\emph{either A or B, but not both}. But TFL \sor is understood in an
\emph{inclusive} way: \emph{A or B, and maybe both}.

Finally, the rules of grammar guarantee that there is no structural ambiguity
in the sentences of TFL. Each sentence can be decomposed to discover exactly
one truth-functional meaning. English, of course, is wildly ambiguous. We will
see a demonstration of this in just a moment.

\subsection{Syntactic Properties of Sentences}

The rules in Def. \ref{def:tfl-sentence} allow us to identify a few important
components of complex sentences.

\begin{definition}[Main Logical Operator]
  \label{def:tfl-main-operator}
  The \emph{main logical operator} of a sentence is the last connective applied
  in its construction.
\end{definition}

Given our definition of TFL sentences, we know that there is a unique main
logical operator for any (non-atomic) sentence. Looked at from the other
direction, we can say that the \emph{scope} of a connective is just the new
sentence that its rule makes.

\begin{definition}[Connective Scope]
  \label{def:tfl-connective-scope}
  The \emph{scope} of a connective is the sentence for which it is the main
  logical operator.
\end{definition}

When checking to see whether some expression is a sentence, we will need to work
backward through these rules: break the expression apart into component parts
until we are left with only atomic sentence letters. Doing this is usually
straight-forward and involves finding the main logical operator (if any) in
the expression under review. Indeed, our bracketing conventions make the main
logical operator and its scope almost immediately apparent. But as we get
better at reading complex expressions, we will treat some brackets more and
more as optional or only occasionally obligatory. For the time being, we will
adopt the conventions in \cite{Mag25} by omitting outermost brackets and
sometimes substituting ``['' and ``]'' for clarity and readability.

\subsection{Parse Trees and Structural Ambiguity}

One strategy that we can use to check the well-formedness of a formula is a
\emph{parse tree}. The idea is that we can \emph{recusrively} break the formula
down based on the grammar rule for the main connective. We can illustrate this
process visually in the form of a tree, but traditionally we use a tree with its
branches pointed down.\footnote{\cite{Smi20} also introduces a notion of
\emph{construction history} (p74) that ``grows'' upward and serves as a
stepping-stone to parse trees. But for our purposes, this seems redundant.}

Consider the following example of a parse tree.

\begin{center}
  \begin{forest}
    for tree={ grow=south, l sep=12pt, s sep=14pt }
    [{$(P \lor \lnot Q) \land (R \to S)$}
      [{$P \lor \lnot Q$}
        [{$P$}]
        [{$\lnot Q$}
          [{$Q$}]
        ]
      ]
      [{$(R \to S)$}
        [{$R$}][{$S$}]
      ]
    ]
  \end{forest}
\end{center}

One of the important properties of parse trees is that each is \emph{uniquely}
associated with a well-formed formula (wff), written at the top (the
\emph{root}). Likewise, every one of its \emph{nodes} and \emph{leaves} is
labeled with a well-formed formula (wff), i.e. a TFL sentence. Of course, this
is just because parse trees are built with a specific process that finds the
main connective and unfolds it by applying the relevant rule in reverse.

\begin{observation}[Parse Tree Uniqueness]
  Every TFL sentence has a unique parse tree.
\end{observation}

Another useful property of parse trees is that we can read important structural
information directly off of them. For example, evey \emph{node} in a parse tree
is decorated with a \emph{subformula} of wff. Indeed, the subformulas are also
exhaustively enumerated by the node decorations.

\begin{observation}[Nodes are subformulas]
  Every node in a parse tree for a sentence \sphi of TFL is decorated with a
  subformula of \ephi.
\end{observation}

Since the \emph{scope} of the main connective of a (sub)formula is just the
(sub)formula itself, we can define this notion via parse trees as well.

\begin{definition}[Scope of a Connective]
  The \emph{scope} of a connective is the subformula decorating the node where
  it is first introduced, reading from the \emph{leaves} of its parse tree up.
\end{definition}

In ordinary English, we can find two distinct types of \emph{ambiguity}.
\emph{Lexical ambiguity} is when on word can mean more than one thing. A popular
example is the word \emph{bank} which can be either a financial institution or
the edge of a river. \emph{Structural ambiguity}, on the other hand, is when a
sentence can have more than one meaning based on how the parts work together. A
popular example is:

\begin{argument}
\item[\textbf{(T)}] The astronomer saw his student with a telescope.
\end{argument}

The sentence in \textbf{(T)} can have two different readings: one where the
astronomer has the telescope and one where the student does. A closely related
concept is \emph{vagueness} which describes terms that are imprecise and often
difficult to apply, such as ``tall'' or ``honorable''. For the purposes of our
logics, we will assume an omniscience about these terms that belies this
vagueness: in our models we will just suppose that we know for each of these
terms exactly when they apply. Like other content terms, their precise meanings
are not important to our rules.

\subsection{Sense and Reference}

Since Frege set out the preliminaries of propositional logic
in the late nineteenth century, logicians have made an important distinction
between two different kinds of meaning that an expression (say a singular term)
can have. On the one hand, it picks out something in the world (or in a domain,
or whatever) that we are taking our sentence to be \emph{about}. Frege called
this type of meaning \emph{reference} (\emph{Bedeutung} in the orginal German).
On the other hand, there is what the expression \emph{means} in some more
abstract way (that is mostly \emph{mental}, but also partly \emph{social}).
This aspect of the meaning of an expression is called \emph{sense}
(\emph{Sinn}).

The touchpoint between our ordinary English content and the symbolization of
TFL sentences is found in the \emph{symbolization key} that we construct for
a given argument. We can think of such a key as a \emph{glossary} of the
content expressions used in the logic. Recall that these expressions in
truth-functional logic are limited to sentences; and that, following Frege, we
take the \emph{reference} of these sentences to be just their \emph{truth
values}. Nevertheless, we can set these values by taking the ``real world''
truth value of the sentences expressed in the symbolization key. In other words,
each entry can be thought of as just a \emph{gloss} of the \emph{sense} of the
content symbols used in the logic. But importantly, what we care about is the
reference that the gloss fixes, not the sense itself. We will see how this works
in a moment.

As intimated above, TFL sentences are \textbf{never} structurally ambiguous.
And we can make efforts in formulating our symbolization keys to minimize, if
not elminate, both vagueness and lexical ambiguity. That said, all three forms
of ambiguity are rampant in ordinary language and our efforts at symbolizing
complex arguments from English will need to proceed with due care. Unlike TFL,
English does not have bracketing rules. We will refer to the structural
ambiguities that can produce more than one TFL sentence as \emph{scope
ambiguities}.

Consider the following sentence and symbolization key:

\begin{argument}
\item[\textbf{(P)}] The professor is literate or not dumb and lazy.
\end{argument}

\begin{grayarg}
\item[A:] The professor is literate.
\item[B:] The professor is dumb.
\item[C:] The professor is lazy.
\end{grayarg}

We can identify a few scope ambiguities, offering plenty of opportunities for
the professor to be illiterate, dumb or lazy but not, as it turns out, all
three. Notice how the bracketing helps to distinguish between the different
\emph{readings} of the English sentence being symbolized.

%\begin{argument}
%  \item[] $(A \eor \enot B) \eand C$
%  \item[] $A \eor \enot (B \eand C)$
%  \item[] $A \eor (\enot B \eand C)$
%\end{argument}

\threeitems{
  $(A \eor \enot B) \eand C$ }{
  $A \eor \enot (B \eand C)$ }{
  $A \eor (\enot B \eand C)$ }

\subsection{Object Language and Metalanguage}

In our informal discussion, it was convenient for us to mix together the
sentences of the logic we were discussing and the sentences about that logic.
Now that we are dealing with formal logic, we will need to be a bit more
careful. We will also need to be clear about when we are using English in an
ordinary way versus when we are using it to talk about a language, whether
English or TFL. To this end, a few more conventions and definitions are in
order.

The most common indicator of level-mixing is quotation marks. Consider the
following examples.

\begin{argument}
\item[A:] ``Aristotle'' is not the original Greek spelling.  
\item[B:] Aristotle taught Alexander the Great.
\end{argument}

Whereas A is about an English word, B is about the ancient philosopher. This is
an important distinction: the former is called \emph{use}; and the latter is
called \emph{mention}. Speaking about the phenomenon more generally, we can
adopt the following definitions.

\begin{definition}[Object Language and Metalanguage]
  When talking about a language such, we call the language that we are talking
  about the \emph{object language} and the language that we are using to talk
  about it the \emph{metalanguage}. 
\end{definition}

For the purpuses of truth-theoretic logic (TFL), our object language is made of
TFL sentences and our metalanguage is English mixed with greek letters and a few
other textual decorations mixed in for precision and convenience.

Since TFL and our metalanguage do not use quotes, we will rarely need to worry
about them. The most important point to make here is that we will use English
as our metalanguage when discussing TFL (and later FOL). But we will make some
adjustments to simplify our presentation. As we have already seen in
\ref{def:tfl-sentence}[Def ], we will use Greek lower case letters ``\ephi'',
``\epsi'',``\echi'', etc as \emph{metavariables} for arbitrary TFL sentences.
Upper case Latin letters ``A'', ``B'', ``C'', etc will continue to be used as
atomic sentence letters. After now, we will rarely mention these variables,
even in the metalanguage. So the quotation marks will almost always be omitted.
\footnote{When they are not, we consider them to be part of the mention.}

One more special feature of our English++ metalanguage is the conclusion marker
(\therefore), which we've seen already in the informal discussion of arguments.
We can use someshorthand to formalize the notation in our metalanguage as well.
For instance, we can describe a sequence of premises and conclusion that
comprise an argument as:

\begin{definition}[Symbolized Argument] Let $\phi_1,\dots,\phi_n$ and \spsi be
  sentences of TFL. $\phi_1,\dots,\phi_n$ $\therefore$ $\epsi$ is a
  \emph{symbolized argument} of TFL. $\phi_1,\dots,\phi_n$ are its
  \emph{premises} and \spsi its conclusion.
\end{definition}

But we will try to strike a healthy balance in the metalanguage between too
much verbiage and too much formalism.

\subsection{Truth-Functionality}

TFL is \emph{truth-functional} as a result of its \emph{semantics}. Put simply,
we evaluate the \emph{truth} of a complex sentences as a \emph{function} of the
truth of their component parts. Ultimately this is a complex function of the
truth of all of its unique atomic sentences. We can define the notion with an
adaptation of the definition in \cite[][p. 72]{Mag25}

\begin{definition}[Truth-Functional]
  \label{def:truth-functional}
  A TFL connective is \emph{truth-functional} iff the truth value of the
  sentence comprising its scope (i.e. the sentence for which it is the main
  logical operator) is uniquely determined by the truth values of the
  sentence(s) it connects.
\end{definition}

Consider our informal definition of \enot\ephi.

\begin{argument}
\item[] A TFL sentence of the form \enot\sphi is true just in case \sphi is
  false. 
\end{argument}

We can see right away how this sort of principle works in practice. The rule is
simple. But actually working with TFL sentences within the system outlined by
TFL requires a substantial shift in how we think about these `sentences'.
Indeed, even this first definition hints at some important features of
(classical) logic that are different than those in ordinary English. We will
discuss these in more detail later, but one is that every TFL sentence is
either true or false; if it is not true, then it is false; if it is not false,
then it is true. This rule (the Law of the Excluded Middle) is exceptionless in
TFL. We will formulate an axiom to that effect later.

\begin{quote}
  PL languages are not mere \emph{shorthand code} for bits of English; their
  whole point is to be usefully tidied-up \emph{replacements} for the messy
  vernacular. (\cite[][p. 83]{Smi20})
\end{quote}

Recall that the process of symbolization for TFL loses a lot of the force or
affect that we often take for granted in ordinary language. This list of
sentences taken from \cite[][p. 73]{Mag25} all have the same truth-value and
can be symbolized with the same TFL sentence ($A \wedge B$), despite a variety
of shaded meanings.

\begin{example}
  Sentences that can be symbolized as $A \wedge B$:
  \begin{enumerate}
  \item Dana is a logician and Dana is a nice person.
  \item Although Dana is a logician, Dana is a nice person.
  \item Dana is a logician despite being a nice person.
  \item Dana is a logician and Dana is a nice person.
  \item Dana is a nice person but also a logician. 
  \item Dana's being a logician notwithstanding, he is a nice person.
  \end{enumerate}
\end{example}

The collapse of these diverse sentences to a single symbolization in TFL happens
precisely \textbf{because} its semantics is truth-functional. To see this, it
will help to look at the important terms one at a time: \emph{truth},
\emph{truth-value} and \emph{function}.

So far we have used ``truth'' to mean something like a metalogical property that
applies to English sentences. On this informal usage, we can fall back on
whatever folk definition we like. This one from American Heritage Dictionary
will do as well as any.

\begin{definition}[Informal Truth] \emph{noun}\footnote{
  The American Heritage® Dictionary of the English Language, 5th Edition}
  \begin{enumerate}
  \item Conformity to fact or actuality. \\ "Does this story have any truth?"
  \item Reality; actuality. \\ "In truth, he was not qualified for the job."
  \item The reality of a situation. \\ "The truth is, she respects your work."
  \end{enumerate}
\end{definition}

We use this notion in understanding our symbolization keys. On our informal
treatment, the following key tells us to consider the truth of A and B to be a
matter of John and Mary's feelings about each other ``in reality.''

\begin{argument}
\item[A:] John loves Mary.
\item[B:] Mary loves John.
\end{argument}

But in TFL, what the key really says is that $A$ has the same \emph{truth value}
that the English sentence ``John loves Mary'' does in reality. We are using
\emph{truth value} here as a technical term.

\begin{definition}[Truth Value]
  A truth value is one of two values: \emph{True} and \emph{False} (symbolized
  in the metalogic as T and F, respectively). Every atomic sentence in TFL is
  assigned one or the other of these values and every complex sentence
  in TFL can be assigned one on the basis of its atomic parts.
\end{definition}

Recall from algebra that a \emph{function} takes one or more inputs
(\emph{arguments}) and produces an output (a \emph{result}). In TFL both the
arguments and the output of truth-functional connectives are truth values. Put
differently, the symbol for a TFL sentence is best understood \textbf{as} its
truth value. The meanings of the English sentences provided in the symbolization
key are just guides to assigning those values properly. Meaning plays no role in
TFL, only truth. This is an important point that takes some time to understand.

\begin{definition}[Truth Function]
  A \emph{truth function} is a function which takes one or more truth values
  as its \emph{argument(s)} and returns a truth value as its output.
\end{definition}

One consequnce of the truth-functional treatment of sentences and connectives is
that we can now fully define the semantics of a TFL connective on the basis of
how the truth of its complex sentence is determined by the truth of the
sentences it connects. We do this with the help of truth tables. Truth-tables
exhaustively spell out the consequences of assigning truth-values to the
components of a sentence. Note that characteristic truth tables use the
metalogical variables \ephi and \epsi because the rules can be applied to
arbitrary TFL sentences, not just to atomic ones.

\subsection{Connectives and their Charateristic Truth Tables}

Recall that we earlier highlighted three generalizations about classical
negation.

%\begin{argument}
%\item[\textbf{LEM:}] \ephi\eor\enot\ephi
%\item[\textbf{LNC:}] \enot(\ephi\eand\enot\ephi)
%\item[\textbf{DN:}]  \enot\enot\ephi\eiff\ephi  
%\end{argument}
\begin{center}
    \textbf{LEM:} \ephi\eor\enot\ephi         \\
    \textbf{LNC:} \enot(\ephi\eand\enot\ephi) \\
    \textbf{DN:} \enot\enot\ephi\eiff\ephi    \\
\end{center}

We can summararize the aggregate of these rules (which, as we will see later,
really are all inter-related) by saying that TFL is \emph{bivalent}.

\begin{definition}[Bivalency]
  A truth-functional language $\mathcal{L}$ is \emph{bivalent} (i.e.
  \emph{two-valued}) iff every sentence \sphi of $\mathcal{L}$ is associated
  with \textbf{exactly one} of the two truth values, T and F, in every
  situation.
\end{definition}

This gives us a straightforward semantics for \enot.

\begin{definition}[Semantics of \enot]
  \label{def:tfl-not}
  For any TFL sentence \ephi, \enot\sphi is true iff \sphi is false.
\end{definition}

We can flesh out the formal definitions for our major connectives by
considering what would happen if we assigned either of the possible truth
values to the sentences that they apply to. We will call this presentation
a \emph{characteristic truth table}. For the negation operator \enot, the
table is simple, because it ``connects'' just one sentence. Let's call it \ephi.

\begin{definition}[Charateristic Truth Table: Negation]
  \label{def:tt-not}
  The semantics of the negation operator \snot is captured by the following
  characteristic truth table: \\
  \begin{center}
    \begin{tabular}{c|c}
      \ephi & \enot\ephi\\
      \hline
      T & F\\
      F & T 
    \end{tabular}
  \end{center}
\end{definition}

There is only one argument: the sentence negated; and that sentence has only
two possible truth values: T and F.

The semantics of \eand is similarly trivial.

\begin{definition}[Semantics of \eand]
  \label{def:tfl-and}
  A TFL sentence of the form \ephi\eand\spsi is true iff both \sphi and \spsi
  are true.
\end{definition}

For the two-place connectives, there are two arguments: one for each of the
sentences connected; and each of those sentence has only two possible truth
values: T and F. If we vary them independently for each sentence, we get four
total possibilities.\footnote{Note that the TTFF and TFTF ordering is
conventional for TFL truth tables. More generally, propositions to the left are
\emph{chunked} and those on the right are \emph{interleaved}. More on this
later.}

\begin{definition}[Charateristic Truth Table: Conjunction]
  \label{def:tt-and}
  The semantics of the connective \sand is captured by the following
  characteristic truth table: \\
  \begin{center}
    \begin{tabular}{c c|c}
      \ephi & \epsi & \ephi\eand\epsi\\
      \hline
      T & T & T\\
      T & F & F\\
      F & T & F\\
      F & F & F
    \end{tabular}
  \end{center}
\end{definition}

Likewise, \sor gets the expected semantics and characteristic truth table.

\begin{definition}[Semantics of \eor]
  \label{def:tfl-or}
  A TFL sentence of the form \ephi\eor\spsi is true just in case either \sphi
  is true or \spsi is true (or both).
\end{definition}

\begin{definition}[Charateristic Truth Table: Disjunction]
  \label{def:tt-or}
  The semantics of the connective \sor is captured by the following
  characteristic truth table: \\
  \begin{center}
    \begin{tabular}{c c|c}
      \ephi & \epsi & \ephi\eor\epsi\\
      \hline
      T & T & T\\
      T & F & T\\
      F & T & T\\
      F & F & F
    \end{tabular}
  \end{center}
\end{definition}

The semantics of conditionals is straight-forward, but not entirely without
controvery. We will present the definition now, but postpone discussion of
the shortcomings.

\begin{definition}[Semantics of \eif]
  \label{def:tfl-if}
  A TFL sentence of the form \ephi\eif\spsi is true just in case either \spsi is
  true or \sphi is false.
\end{definition}  

\medskip

\begin{definition}[Charateristic Truth Table: Conditional]
  \label{def:tt-if}
  The semantics of the connective \sif is captured by the following
  characteristic truth table: \\
  \begin{center}
    \begin{tabular}{c c|c}
      \ephi & \epsi & \ephi\eif\epsi\\
      \hline
      T & T & T\\
      T & F & F\\
      F & T & T\\
      F & F & T
    \end{tabular}
  \end{center}
\end{definition}

\medskip

\begin{definition}[Semantics of \eiff]
  \label{def:tfl-iff}
  A TFL sentence of the form \ephi\eiff\sphi is true just in case \sphi and
  \spsi are either both true or both false.
\end{definition}

\medskip

\begin{definition}[Charateristic Truth Table: Biconditional]
  \label{def:tt-iff}
  A TFL sentence of the form \ephi\eiff\sphi is true just in case \sphi and
  \spsi are either both true or both false. The semantics of the connective
  `\eiff' is defined by the following characteristic truth table: \\
  \begin{center}
    \begin{tabular}{c c|c}
      \ephi & \epsi & \ephi\eiff\epsi\\
      \hline
      T & T & T\\
      T & F & F\\
      F & T & F\\
      F & F & T
    \end{tabular}
  \end{center}
\end{definition}

Instead of depending on our knowledge of reality to assign truth-values to
atomic sentencs as we do in the informal assessment of truth in our
symbolization tables, we can simply assign values to each of the atomic
sentence letters used in our logic. Recall that we think of a symbolization
key as a glossary of the \emph{senses} of the atomic sentences in our logic. But
TFL expects their \emph{references} (i.e. the truth values themselves). We will
define a \emph{valuation} as the specification of the truth values for all the
atomic sentences in a given instance of TFL.

\subsection{TFL Valuations}

\begin{definition}[Valuation]
  \label{def:valuation}
  A \emph{valuation} is any assignment of truth values to the atomic sentence
  letters of TFL. On a valuation, each atomic sentence letter is assigned
  exaclty one truth value.
\end{definition}

The completeness of a TFL valutation is refered to as \emph{determinacy}.

\begin{definition}[Determinacy]
  For any TFL language, once we have stipulated the senses of its atoms--fixed
  their truth conditions--each of these atoms will be determinately true or
  determinately false in any possible situation. The atoms of a TFL language
  are to be interpreted so that they have determinate truth values, exactly
  one of T and F, in any situation. (modified from \cite[][p. 84-5]{Smi20})
\end{definition}

But, of course, ordinary English often lacks such determinacy. Nevertheless,
we will see that it is an important aspect of classical logical systems. One
such aspect is the resulting determinacy of \textbf{all} the sentences of TFL,
both atomic and complex.

% valuations uniquely determine every wff (p89)
\begin{observation}[Valuations uniquely determine all truth]
  \label{obs:unique-determinacy}
  Note that determinacy is contagious via rule application.
  
  \begin{enumerate}
  \item A valuation of all the atomic sentences of any TFL sentence \sphi
    \emph{uniquely determines} a resulting valuation of \ephi.
  \item A valuation of all the atomic sentences of any TFL language
    $\mathcal{L}$ \emph{uniquely determines} a resulting valuation of
    $\mathcal{L}$, however complex.
  \end{enumerate}
  
\end{observation}

\subsection{Setting up Simple Truth Tables}

% mini-tables for deterministic valuations (pp90-2)
We can take a deterministic valuation to be a stand-in for what we were earlier
calling a \emph{situation}. With this meaning, a situation would just be the
state of affairs defined by the truth values as they are set in the valuation.
A truth table for such a valuation is easy to construct because we only need to
consider one option for each sentence: the actual truth value that is assigned
to it under the valuation. It will have just one row.

\begin{center}
  \begin{tabular}{c c|c c c c c}
    $A$ & $B$ & $(A$ & \eand & $B)$ & \eif       & $A$ \\
    \hline
    F   & T   & F    & F     & T    & \textbf{T} & F
  \end{tabular}
\end{center}

The truth value (F) of the entire formula ($(A \eand B) \eif A$) is bolded
in the column under the main connective. This formula is \emph{vacuously true},
because its antecedent is false.

But full truth tables are \emph{exhaustive} in the sense that every possible
valuation for a given set of sentences can be represented as a row in the
table, with $2^n$ such rows for a TFL sentence that has $n$ different atomic
sentence letters. How this works in practice is best demonstrated by example.
Consider the TFL sentence $(A\wedge B) \rightarrow A$. We can see how the
truth of the whole depends on the truth of its parts by working through each of
the possible valuations ``from the inside out''.

First, set up the table with all of the possible atomic truth values.

\begin{center}
  \begin{tabular}{c c|c c c c c}
    $A$ & $B$ & $(A$ & \eand & $B)$ & \eif & $A$ \\
    \hline
    T & T & & & & &\\
    T & F & & & & &\\
    F & T & & & & &\\
    F & F & & & & &
  \end{tabular}
\end{center}

Then add the atomic truth values to the appropriate columns on the right side of
the table.

\begin{center}
  \begin{tabular}{c c|c c c c c}
    $A$ & $B$ & $(A$ & \eand & $B)$ & \eif & $A$ \\
    \hline
    T & T & T & & T & & T\\
    T & F & T & & F & & T\\
    F & T & F & & T & & F\\
    F & F & F & & F & & F
  \end{tabular}
\end{center}

Now we can evaluate the truth of the sentences with connectives, starting with
the ``innermost'' ones (i.e. the ones that connect only atomic sentences). Here
we start with the \sand connective.

\begin{center}
  \begin{tabular}{c c|c c c c c}
    $A$ & $B$ & $(A$ & \eand & $B)$ & \eif & $A$ \\
    \hline
    T & T & T & \bf{T} & T & & T\\
    T & F & T & \bf{F} & F & & T\\
    F & T & F & \bf{F} & T & & F\\
    F & F & F & \bf{F} & F & & F
  \end{tabular}
\end{center}

We continue the process until we reach the main logical operator (i.e. the
outermost connective). The truth value assigned at this column is the truth
value for the entire sentence.

\begin{center}
  \begin{tabular}{c c|c c c c c}
    $A$ & $B$ & $(A$ & \eand & $B)$ & \eif & $A$   \\
        &     &      & \ephi &      & \eif & \epsi \\
    \hline
    T & T & T & T & T & \bf{T} & T\\
    T & F & T & F & F & \bf{T} & T\\
    F & T & F & F & T & \bf{T} & F\\
    F & F & F & F & F & \bf{T} & F
  \end{tabular}
\end{center}

Notice how we can use the metalogical sentence variables as substitutes for more
complex TFL sentences. Here they are really being used as \emph{logical}
sentence variables, representing not arbitrary sentences but these specific two:

\begin{argument}
  \item[\ephi:] $A \wedge B$
  \item[\epsi:] $A$
\end{argument}

Returning to our example from last week about the literacy of a certain
professor, we can see how this approach becomes complicated fast. Recall that
the three readings we found for our example were:

\begin{argument}
  \item[\ephi:] $(A \eor \enot B) \eand C$
  \item[\epsi:] $A \eor \enot (B \eand C)$
  \item[\echi:] $A \eor (\enot B \eand C)$
\end{argument}

We can explore the different truth conditions for these sentences by expressing
each of them in the form of a truth table. Now that we have three atomic
sentences ($A$, $B$, and $C$), we will need $2^3 = 8$ rows to express all of
the possible valuations relevant to this argument. These valuations are
straight-forward, if time-consuming, to lay out no matter how many atomic
sentences occur in our argument.

Now ``possible'' can mean a few different things, depending on how it is
interpreted. Consider the following symbolization key.

\begin{argument}
\item[A:] John has his cake.
\item[B:] John ate his cake.
\end{argument}

In one sense, we can just list out all of this possibilities. In that case,
$A \eand B$ will be true when both A and B are true. We can say that it is
\emph{c-possible} that both A is true and B is true, since this is one of the
possibile combinations of true and false required by the truth-table method.

On the other hand, we all know that you can't $A \eand B$. Once you've eaten
the cake (B), it is gone (\enot A). Although treating both A and B as true is
combinatorily possible, the situtation that it describes is not
\emph{realizable} and so we can say that it is not \emph{r-possible}.

We will not lean heavily on these definitions except to note that when we talk
about possible situations in TFL, we mean that they are c-possible.

\subsection{Building Complete Truth Tables}

\begin{definition}[Complete Truth Table]
  The complete truth table for a TFL sentence \ephi is a truth table with each
  c-possible valuation; that is, for each possible combination of truth value
  assignments to all the unique atomic sentences found in \ephi.

  We can construct a complete truth table for a TFL sentence with $n$ unique
  atomic sentences using the following procedure:
  
  \begin{enumerate}
  \item Write each atomic sentence letter and the TFL sentence across the top,
    assuming a column for each of its connectives and atomic sentences.
  \item Insert $2^n$ rows.
  \item Populate the first column (the first $n$ columns are the unique atomic
    sentence letters) by adding ``T'' to the first half ($2^n/2$) of the rows
    and then ``F'' for the remaining half.
  \item For each subsequent column $i<=n$, fill in the portion of the column
    directly to the right of each uninterrupted series of ``T''s (or ``F''s) by
    filling in the first half with ``T'' and the second half with ``F''.
  \item If done correctly, the last column of atomic letters (column $n$) should
    alternate ``T'' then ``F'' a total of $2^n/2$ times.
  \item Fill in the remaining columns according to the parse tree and connective
    semantics as above.
  \end{enumerate}

\end{definition}

Now that we know how to build a truth table for three atomic sentences, we can
build truth tables knowing that they are complete. First, set up the table with
all of the possible atomic truth values.

%$A$ & $B$ & $C$ & $A$  & \eor & \enot    & $(B$ & \eand & $C)$ \\
%$A$ & $B$ & $C$ & $(A$ & \eor & $(\enot$ & $B$  & \eand & $C$  \\
\begin{center}
  \begin{tabular}{c c c|c c c c c c}
    $A$ & $B$ & $C$ & $(A$ & \eor & \enot&$B)$ & \eand & $C$  \\
    \hline
    T & T & T & & & & & &\\
    T & T & F & & & & & &\\
    T & F & T & & & & & &\\
    T & F & F & & & & & &\\    
    F & T & T & & & & & &\\
    F & T & F & & & & & &\\
    F & F & T & & & & & &\\
    F & F & F & & & & & &
  \end{tabular}
\end{center}

Copy the truth values of the atomic sentences into the approriate columns on the
right side of the table.

\begin{center}
  \begin{tabular}{c c c|c c c c c c}
    $A$ & $B$ & $C$ & $(A$ & \eor & \enot&$B)$ & \eand & $C$  \\
    \hline
    T & T & T & T & & & T & & T\\
    T & T & F & T & & & T & & F\\
    T & F & T & T & & & F & & T\\
    T & F & F & T & & & F & & F\\    
    F & T & T & F & & & T & & T\\
    F & T & F & F & & & T & & F\\
    F & F & T & F & & & F & & T\\
    F & F & F & F & & & F & & F
  \end{tabular}
\end{center}

Apply the connective rules up the parse tree, starting with the atomic
sentences.\footnote{Starting here, we will use the convention from \cite{Smi20}
where multiple new steps are indicated in the same table by numbering along the
bottom.}

\begin{center}
  \begin{tabular}{c c c|c c c c c c}
    $A$ & $B$ & $C$ & $(A$ & \eor & \enot&$B)$ & \eand & $C$  \\
    \hline
    T & T & T & T & T & F & T & \bf{T} & T\\
    T & T & F & T & T & F & T & \bf{F} & F\\ 
    T & F & T & T & T & T & F & \bf{T} & T\\
    T & F & F & T & T & T & F & \bf{F} & F\\    
    F & T & T & F & F & F & T & \bf{F} & T\\
    F & T & F & F & F & F & T & \bf{F} & F\\
    F & F & T & F & T & T & F & \bf{T} & T\\
    F & F & F & F & T & T & F & \bf{F} & F\\
    \hline
      &   &   &   & 2 & 1 &   & 3      &  \\
  \end{tabular}
\end{center}

The truth of the main sentence is provided for each valuation in the bolded
column. Similarly, the truth conditions for the other two sentences are
expressed in the bolded columns of the following two complete truth tables.

\begin{center}
  \begin{tabular}{c c c|c c c c c c}
    $A$ & $B$ & $C$ & $A$ & \eor & \enot & $(B$ & \eand & $C)$ \\
    \hline
    T & T & T & T & \bf{T} & F & T & T & T\\
    T & T & F & T & \bf{T} & T & T & F & F\\
    T & F & T & T & \bf{T} & T & F & F & T\\
    T & F & F & T & \bf{T} & T & F & F & F\\    
    F & T & T & F & \bf{F} & F & T & T & T\\
    F & T & F & F & \bf{T} & T & T & F & F\\
    F & F & T & F & \bf{T} & T & F & F & T\\
    F & F & F & F & \bf{T} & T & F & F & F\\
    \hline
      &   &   & 1 & \bf{6} & 5 & 2 & 4 & 3
  \end{tabular}
\end{center}

\begin{center}
  \begin{tabular}{c c c|c c c c c c}
    $A$ & $B$ & $C$ & $(A$ & \eor & $(\enot$ & $B$  & \eand & $C$  \\
    \hline
    T & T & T & T & \bf{T} & F & T & F & T\\ 
    T & T & F & T & \bf{T} & F & T & F & F\\
    T & F & T & T & \bf{T} & T & F & T & T\\
    T & F & F & T & \bf{T} & T & F & F & F\\    
    F & T & T & F & \bf{T} & F & T & F & T\\
    F & T & F & F & \bf{F} & F & T & F & F\\
    F & F & T & F & \bf{T} & T & F & T & T\\
    F & F & F & F & \bf{F} & T & F & F & F\\
    \hline
      &   &   & 1 & \bf{6} & 4 & 2 & 5 & 3
  \end{tabular}
\end{center}

If we compare the results in a single collapsed table, it is easy to see that
these TFL sentences have different truth conditions; that is, these sentences
not true at the same valuations. I've bolded the outliers to emphasize the
differences.

\begin{center}
  \begin{tabular}{c c c|c c c}
    $A$ & $B$ & $C$ & \ephi & \epsi & \echi \\
    \hline
    T & T & T &     T  &     T   &     T \\
    T & T & F & \bf{F} &     T   &     T \\
    T & F & T &     T  &     T   &     T \\
    T & F & F & \bf{F} &     T   &     T \\    
    F & T & T &     F  &     F   & \bf{T}\\
    F & T & F &     F  & \bf{T}  &     F \\
    F & F & T &     T  &     T   &     T \\
    F & F & F &     F  & \bf{T}  &     F
  \end{tabular}
\end{center}

Next lecture, we will discuss the semantic implications of these differences in
detail. And we will look at some ways to simplify our process of constructing
truth tables.

\subsection{Expressive Adequacy}

Armed with truth tables and a new found confidence in working through them,
we can begin to explore some fundamental questions about the logic itself.
Are these all of the truth-functional connectives that can be defined? Do we
need them all or will just a few suffice? These are important topics in
\emph{metalogic}, the major focus of Logic 2. But we shouldn't be afraid to
think through them a little here, using our new toolbox.

First, notice that we don't really need both \sand \textbf{and} \eor.

% DeMorgan (p113)
\begin{observation}[\sand and \sor are mutually-definable with \enot]
  The connectives \sand and \sor could each be used in a definition of the
  other.
  \begin{enumerate}
  \item The sentence \ephi\eand\spsi can be expressed without \sand as
    \enot(\enot\ephi\eor\enot\epsi).
 \item The sentence \ephi\eor\spsi can be expressed without \sor as
    \enot(\enot\ephi\eand\enot\epsi).  
  \end{enumerate}
  
\end{observation}

We can see this by comparing them in truth tables.

\twoitems{
  \begin{center}
    \begin{tabular}{c c||c|c c c c}
      \ephi&\epsi&\ephi\eand\epsi& \enot    &$(\enot\ephi$&\eor&$\enot\epsi)$\\
      \hline
        T  &  T  &\textbf{T}     &\textbf{T}&     F       &  F &  F          \\
        T  &  F  &\textbf{F}     &\textbf{F}&     F       &  F &  T          \\
        F  &  T  &\textbf{F}     &\textbf{F}&     T       &  F &  F          \\
        F  &  F  &\textbf{F}     &\textbf{F}&     T       &  T &  T          \\
    \end{tabular}
  \end{center}
}{
  \begin{center}
    \begin{tabular}{c c||c|c c c c}
      \ephi&\epsi&\ephi\eor\epsi& \enot    &$(\enot\ephi$&\eand&$\enot\epsi)$\\
      \hline
        T  &  T  &\textbf{T}     &\textbf{T}&     F       &  F &  F          \\
        T  &  F  &\textbf{T}     &\textbf{T}&     F       &  F &  T          \\
        F  &  T  &\textbf{T}     &\textbf{T}&     T       &  F &  F          \\
        F  &  F  &\textbf{F}     &\textbf{F}&     T       &  T &  T          \\
    \end{tabular}
  \end{center}
}

Similarly, we can dispense with \eif as long as we have \eor and \enot.

\begin{observation}[\sif can be defined with \sor and \enot]
  The sentence \ephi\eif\spsi can be expressed without \sif as
  \enot\ephi\eor\epsi.
\end{observation}

\begin{center}
  \begin{tabular}{c c||c|c c c}
    \ephi&\epsi&\ephi\eif\epsi&\enot\ephi&\eor&\epsi \\
    \hline
      T  &  T  &\textbf{T}    &    F     &\textbf{T}&T \\
      T  &  F  &\textbf{F}    &    F     &\textbf{F}&F \\
      F  &  T  &\textbf{T}    &    T     &\textbf{T}&T \\
      F  &  F  &\textbf{T}    &    T     &\textbf{T}&F
  \end{tabular}
\end{center}

Finally, we can dispense with \siff as long as we have \sand and \eif.

% A<->B is (A->B)&(B->A) (p 157)
\begin{observation}[\siff can be defined with \sif and \eand]
  The sentence \ephi\eiff\spsi can be expressed without \siff as
  $(\ephi\eif\epsi)\eand(\epsi\eif\ephi)$. Note that we can further reduce
  this, eliminating even \sand and \sif as before.
\end{observation}

\begin{center}
  \begin{tabular}{c c||c|c c c}
    \ephi&\epsi&\ephi\eiff\epsi&\ephi\eif\epsi&   \eand    &\epsi\eif\ephi \\
    \hline
      T  &  T  &   \textbf{T}  &      T       &\textbf{T}  &       T \\
      T  &  F  &   \textbf{F}  &      F       &\textbf{F}  &       T \\
      F  &  T  &   \textbf{F}  &      T       &\textbf{F}  &       F \\
      F  &  F  &   \textbf{T}  &      T       &\textbf{T}  &       T \\
  \end{tabular}
\end{center}

Taken together, these results suggest that we could do without the connectives
\eand, \eif, and \eiff, building our system entirely from \sor and \snot.

\subsubsection{Are there more connectives?}

Before seeing just how far we can strip our system down, let's take a look at
the first of our questions: is this all of the truth-functional connectives that
could be defined? It seems the answer is: sort of.

% XOR with existing connectives (p 114)
We can explore a few well-known candidates to get a sense of what I mean here.
First, consider the alternative style of disjunction common in ordinary
English, \emph{exclusive disjunction} (XOR), which means ``either A or B, but
not both''. Let's suppose that XOR is symbolized in TFL as a connective
$\oplus$. It's easy to see that such a connective is definable using \sand,
\sor and \enot.

\begin{center}
  \begin{tabular}{c c||c|c c c c c}
    $\ephi$&$\epsi$&$\ephi\oplus\epsi$&$(\ephi\eor\epsi)$&\eand     &\enot&$(\ephi\eand\epsi)$ \\
    \hline
      T    &  T    &   \textbf{F}      &         T        &\textbf{F}& F   &          T        \\
      T    &  F    &   \textbf{T}      &         T        &\textbf{T}& T   &          F        \\
      F    &  T    &   \textbf{T}      &         T        &\textbf{T}& T   &          F        \\
      F    &  F    &   \textbf{T}      &         F        &\textbf{T}& T   &          F        \\
  \end{tabular}
\end{center}

\begin{definition}[Expressively Adequate]
  A collection of truth-functional connectives is \emph{expresively adequate}
  iff they can be used collectively to define any other possible connective.
\end{definition}

% $(x,y,z) with existing connectives (pp 115-6)
%\begin{definition}[Charateristic Truth Table: $\dagger$]
%  \label{def:tt-dagger}
%  The semantics of the connective $\dagger$ is captured by the following characteristic truth table: \\
%  \begin{center}
%    \begin{tabular}{c c c|c}
%      \ephi&\epsi&\echi& $\dagger(\ephi,\epsi,\echi)$ \\
%      \hline
%      T    & T   & T   &  F
%      T    & T   & F   &  T
%      T    & F   & T   &  F
%      T    & F   & F   &  T
%      F    & T   & T   &  F
%      F    & T   & F   &  F
%      F    & F   & T   &  T
%      F    & F   & F   &  F
%    \end{tabular}
%  \end{center}
%\end{definition}

\begin{observation}[TFL is Expressively Adequate]
  We can define any new possible n-place truth-functional connective $\dagger$
  given only the connectives already available in TFL: \eand, \sor and
  \enot.\footnote{See \cite[][pp. 115-117]{Smi20} for a more complete
  justification of this observation.}
\end{observation}

\subsubsection{How few connectives are adequate?}

So we don't need any more connectives. But how few can we get away with? We
already know that just \enot and \eor is expressively adequate for TFL. As is
\snot and \eand. So two? Can we get away with one? The answer is yes, but we
need to make a new one. Actually, there are two options.

\begin{definition}[Charateristic Truth Tables: Peirce's Arrow and the Sheffer stroke]
  \label{def:tt-downarrow}
  The semantics of the connectives $\downarrow$ and $\uparrow$ are captured by
  the following characteristic truth tables: \\
  \twoitems{
    \begin{center}
      \begin{tabular}{c c|c}
        \ephi & \epsi & $\ephi\downarrow\epsi$\\
        \hline
        T & T & F\\
        T & F & F\\
        F & T & F\\
        F & F & T
      \end{tabular}
    \end{center}
  }{
    \begin{center}
      \begin{tabular}{c c|c}
        \ephi & \epsi & $\ephi\uparrow\epsi$\\
        \hline
        T & T & F\\
        T & F & T\\
        F & T & T\\
        F & F & T
      \end{tabular}
    \end{center}
  }
\end{definition}

Truth tables for the following equivalences are left as an exercise.

\begin{center}
  \begin{tabular}{c|c|c|c}
    &
    TFL &
    Peirce's Arrow ($\downarrow$) &
    Sheffer Stroke ($\uparrow$) \\
    \hline
    Negation &
    \enot\ephi &
    $(P\downarrow P)$ &
    $(P\uparrow P)$ \\
    Disjunction &
    \ephi\eor\epsi &
    $(P \downarrow Q)\downarrow(P \downarrow Q)$ &
    $(P \uparrow P)\uparrow(Q \uparrow Q)$ \\
    Conjunction &
    \ephi\eand\epsi &
    $(P \downarrow P)\downarrow(Q \downarrow Q)$ &
    $(P \uparrow Q)\uparrow(P \uparrow Q)$ \\
  \end{tabular}
\end{center}

